% mage-auto-heading: engineering
% mage-auto-heading: software-agents
% mage-auto-heading: mage
% Authoritative LaTeX manuscript; Markdown is a generated compatibility projection.

% <!-- part-title: Orientation and Reference -->

% <!-- chapter-title: Glossary -->

% <!-- noqa: book-section-cap — glossary; each group is an intentional scannable list of terms under one heading -->

\chapter{Glossary}

This glossary groups the book’s vocabulary under Engineering, Software Agents, and MAGE.

\section{Engineering}\label{heading:engineering}

\textbf{Authority.} Who is entitled to exercise engineering judgment: to decide what a system must do, permit or refuse an action, and answer for the result? Authority is a matter of decision rights and responsibility, held by a person or institution and delegated deliberately. It is distinct from \emph{enforcement}---the mechanisms through which the environment holds work to an obligation---and from an \emph{authoritative source}, the canonical origin of a fact.

\textbf{Engineering.} The discipline of exercising informed \emph{control} over consequential systems and accepting \emph{responsibility} for their outcomes. \emph{Control} means retaining sufficient command of a system to understand and direct it, evaluate the evidence for its consequential properties, recognize when its assumptions fail, and intervene when necessary. \emph{Responsibility} means remaining answerable for the consequential decisions made under one\textquotesingle s authority and for the resulting system. Engineering does not require personally performing every act of implementation, and work may be delegated to other humans or to machine intelligence. Engineering is defined by the control, evidence, expertise, and responsibility exercised, regardless of whether the work belongs to a conventionally named engineering discipline. Software engineering applies this discipline to software systems.

\textbf{Engineering judgment.} The use of technical knowledge, evidence, and experience to decide what matters in an engineering situation and what action is justified. It includes choosing appropriate models and assumptions, assessing whether evidence is sufficient, and weighing the consequences for people affected by the system~\citep{petroski1994judgment, swenson2025judgment}. Competence in making these decisions, together with accountability for their consequences, helps justify society entrusting engineers with authority over consequential systems~\citep{engineeringcouncil2026ethics}.

\textbf{Engineering obligation.} A requirement or engineering decision that the responsible engineering authority has determined a system or its development must satisfy. Obligations may arise externally---from users, regulations, standards, contracts, or operating environments---or internally, from decisions about architecture, behavior, quality, or process. An obligation is evaluable when the environment can obtain evidence about whether it holds. It is enforceable when that evidence can be connected to a mechanism that prevents or refuses violations within a declared scope. Enforcement is the application of that mechanism.

\textbf{Model.} A purposeful simplification of a system that preserves what matters for a particular engineering question. By leaving other details out, a model makes the relevant properties easier to understand and reason about than the system itself.

\textbf{Software engineering.} The engineering discipline associated with the creation and modification of software systems. Due to the flexibility of the engineering medium, software engineering is unusual in the rate of change of the engineered system. 

% <!-- index-def: tolerance -->

% <!-- index-def: degrees-of-freedom -->

\textbf{Tolerance and degree of freedom.} Engineering obligations need not determine a single acceptable implementation. Tolerances describe the bounds of acceptable variation; degrees of freedom describe the choices remaining within those bounds. Note that an unmodeled choice may not truly be free. If some alternatives would violate an unwritten engineering obligation, the choice remains constrained even if the model does not express that constraint.

\textbf{Traceability.} Explicit, machine-readable links from model elements to the artifacts, evidence, obligations, or other model elements they represent, constrain, or correspond to. A link can remain intact while its meaning has gone stale: a test may still carry a requirement's identifier after it has stopped testing that requirement.

\section{Software Agents}\label{heading:software-agents}

\textbf{Agent.} An AI system that pursues an open-ended or underspecified task through multiple steps, exercising judgment about what to do next and using tools in response to intermediate results. This book is focused specifically on the use of ``coding agents'' in a software development context. 

\textbf{Agentic harness.} The runtime that turns a foundation model’s reasoning into action by managing context, invoking tools, and returning results to the model. It may expose hooks at which additional code can run. How those facilities are configured and used is part of the engineered environment.

% <!-- index-def: churn -->

\textbf{Churn.} Work spent rediscovering context, undoing recent changes, repairing regressions, or reconciling inconsistencies instead of advancing the system.

% <!-- index-def: commodity-intelligence -->

\textbf{Commodity intelligence.} General-purpose machine intelligence available cheaply and at enough scale that implementation becomes less scarce relative to engineering judgment. The term describes an economic change, not a claim that such intelligence is free, unlimited, excellent, or sufficient for engineering.

\textbf{Foundation model.} The trained model at the bottom of the agent stack. Its capabilities and context capacity shape the harness and engineered environment around it.

\textbf{Reasoning.} The use of available information to draw inferences, evaluate alternatives, and work toward a solution. A reasoning horizon is the practical limit of the work a reasoner and its supporting machinery can carry through reliably, given the available representations and resources. For an agent, this depends on the foundation model together with the harness and engineered environment.

\textbf{Scaffolding.} Additional guidance or structure that helps a particular implementer satisfy an engineering obligation without changing the obligation itself. More capable implementers may require less scaffolding to satisfy the same obligation.

\textbf{Skills and tools.} Capabilities exposed through an agentic harness that extend what an agent can do. A skill provides reusable guidance for a recurring class of work; a tool provides a capability the agent can invoke rather than approximate through reasoning.

\section{MAGE}\label{heading:mage}

\textbf{Agentic engineering.} The engineering of systems and environments in which autonomous agents perform consequential work. MAGE develops one approach, using software engineering as its principal example.

\textbf{Alignment (MAGE).} The engineering activity of making selected obligations enforceable through constraints, evidence, and admission mechanisms in a governed engineering environment. Unlike model--implementation correspondence, Alignment is not limited to establishing agreement between representations. It does not require changing the foundation model's objectives or behavior.

Active alignment changes a system to satisfy modeled obligations. Passive alignment protects those obligations during subsequent engineering. The alignment kernel comprises the models and controls that protect properties an upcoming engineering transformation must preserve. Engineers can extend or revise the kernel as the system and its obligations change. Chapters \magevalue{chapter:alignment|num} and \magevalue{chapter:one-problem-many-models|num} develop these ideas.

% <!-- index-def: compounding -->

\textbf{Engineering capital.} Durable engineering structure---models, packaged procedures, mechanisms, canonical documentation, and other reusable assets---that makes later work cheaper or less uncertain. Governance conversion creates such structure; it becomes capital when later work benefits from it, and depreciates when it no longer fits the system or costs more to maintain than it returns.

\textbf{Governance.} Deciding what autonomous work must know, what it may do, what evidence it must produce, and what mechanisms will hold those decisions over time.

\textbf{Governance conversion.} Turning a lesson from a recurring failure, repeated judgment, or observed gap into durable engineering structure that future work inherits. The structure may be an engineering model, packaged procedure, or mechanism that enforces an obligation. The conversion lets later work use the resulting structure without reconstructing the original judgment.

\textbf{Governance mechanisms.} Repeatable structures in the engineered environment that constrain agentic work, produce evidence about it, or control whether it may proceed. This book distinguishes three kinds of mechanisms.

\begin{itemize}
\item
  \textit{Constraint.} Restricts available actions or states.
\item
  \textit{Sensor.} Observes the work or system and produces evidence about artifacts, proposed actions, or executions.
\item
  \textit{Gate.} Evaluates whether an action or result may proceed and enforces the decision.
\end{itemize}

\textbf{Governed Engineering Environment.} The engineered environment surrounding autonomous work, comprising the existing system and the representations, tools, processes, permissions, and controls through which engineers direct and govern its development. Repeatable mechanisms enforce selected engineering obligations, allowing consequential decisions to persist across successive agents and changes.

\textbf{Model drift.} Disagreement between a model and the system, evidence, or other representation to which it is meant to correspond. The disagreement identifies a correspondence problem whose resolution may require engineering judgment, especially if it is unclear which view is authoritative. Reconciliation often requires aligning the system with the engineering intent captured in the model. See also: \textit{Ratcheteering}.

\textbf{Modeling Principle.} Externalize knowledge and intent into engineering models that both engineers and agents can reason through. Appropriate representations make otherwise intractable engineering questions manageable and make additional properties available for Alignment.

% <!-- index-def: ratcheteering -->

\textbf{Ratcheteering.} The practice of introducing an authoritative engineering model that deliberately disagrees with an existing system, then directing agents to change the implementation until it conforms to the model. Rather than adapting the model to describe the existing implementation, the engineer uses the model to establish how the system should be structured or behave. See also: \textit{Model drift}.

\textbf{Semantic gap.} The distance between where an engineering obligation arises and where the environment has enough meaning and evidence to decide whether the obligation is satisfied. The gap explains why some checks must wait until a later boundary, when enough of the system has been assembled to make the obligation decidable.

\textbf{Software factory.} The engineered system through which software changes are realized from engineering intent and evaluated, admitted, deployed, and improved. Whereas the Governed Engineering Environment emphasizes how autonomous work is governed, the software factory emphasizes the organization of production. \magevalue{chapter:software-factories} develops this perspective.
